\documentclass[a4paper,11pt]{article}
\usepackage{amsmath,amssymb,amsthm}
\usepackage[margin=2.5cm]{geometry}
\usepackage{bm}
\usepackage{algorithm}
\usepackage{algpseudocode}
\usepackage{booktabs}

\newcommand{\va}{\bm{a}}
\newcommand{\vb}{\bm{b}}
\newcommand{\vc}{\bm{c}}
\newcommand{\vct}{\tilde{\bm{c}}}
\newcommand{\ve}{\bm{e}}
\newcommand{\vu}{\bm{u}}
\newcommand{\vv}{\bm{v}}
\newcommand{\Xb}[2]{X_{#1}^{#2}}
\newcommand{\Xn}[2]{\bar{X}_{#1}^{#2}}
\newcommand{\Gb}[1]{G_{#1}}
\newcommand{\frob}[1]{\lVert #1 \rVert_{\mathrm{F}}}
\newcommand{\tr}{\mathrm{tr}}
\newcommand{\Nk}{N_k}
\newcommand{\Nt}{N_\tau}

\theoremstyle{plain}
\newtheorem{theorem}{Theorem}
\newtheorem{remark}{Remark}

\title{Matching Solution Method for Joint Eigenfunctions\\
       of Cosmological Perturbation Bases}
\author{Will Handley \and Wei Ning Deng}
\date{\today}

\begin{document}
\maketitle

\begin{abstract}
We address the problem of finding coefficient vectors $\va, \vb \in \mathbb{R}^{\Nk}$
such that the time evolutions $\Xb{1}{p}\va$ and $\Xb{2}{p}\vb$ are \emph{equal} for
every perturbation type $p$ simultaneously, not merely correlated.
We show that canonical correlation analysis (CCA, see \texttt{cca\_formulation.pdf})
is insufficient for this goal because it maximises a single aggregate correlation,
allowing strong agreement in some fields to mask a poor match in others.
The correct formulation is to minimise the sum of squared per-field residuals, which
reduces to the minimum-eigenvalue problem on a $(2\Nk\times 2\Nk)$ block matrix.
The framework is equivalent to finding a vector in the intersection of null spaces,
a classical construction in linear algebra.
\end{abstract}

%----------------------------------------------------------------------
\section{Problem Statement}
\label{sec:problem}
%----------------------------------------------------------------------

The setup is the same as in \texttt{cca\_formulation.pdf}.
A Boltzmann code produces solutions for perturbation variables
$\delta_r, \delta_m, v_r, v_m$ as functions of conformal time $\tau$ and
wavenumber $k$.
For $\Nk$ wavenumbers we have two bases:
\begin{itemize}
    \item \textbf{Basis 1} (integer-$k$): $\Xb{1}{p} \in \mathbb{R}^{\Nt\times\Nk}$,
    \item \textbf{Basis 2} (allowed-$k$): $\Xb{2}{p} \in \mathbb{R}^{\Nt\times\Nk}$,
\end{itemize}
for each perturbation type $p$.

\textbf{Goal.} Find sparse coefficient vectors $\va, \vb \in \mathbb{R}^{\Nk}$
such that
\begin{equation}
    \Xb{1}{p}\va = \Xb{2}{p}\vb \qquad \text{for all } p.
    \label{eq:goal}
\end{equation}
This is a stricter requirement than the CCA goal of maximising the correlation
between $\Xb{1}{p}\va$ and $\Xb{2}{p}\vb$: we want the two time-series vectors
to be \emph{identical}, not merely proportional.

%----------------------------------------------------------------------
\section{Limitation of CCA for Exact Matching}
\label{sec:cca-limit}
%----------------------------------------------------------------------

The CCA objective is
\begin{equation}
    \rho(\va,\vb) = \frac{\va^\top \Gb{12}\,\vb}{\sqrt{\va^\top \Gb{1}\,\va
    \;\cdot\; \vb^\top \Gb{2}\,\vb}},
    \label{eq:cca}
\end{equation}
where $\Gb{12} = \sum_p w_p (\Xn{1}{p})^\top \Xn{2}{p}$ aggregates the
cross-field agreement into a \emph{single scalar} via a weighted sum over $p$.

Two structural problems prevent CCA from guaranteeing exact per-field matching.

\paragraph{Problem 1: outvoting.}
A high value $\rho \approx 1$ only means that the \emph{weighted average}
agreement across all fields is high.
If $\delta_r, v_r, v_m$ match well but $\delta_m$ does not, the three
well-matching fields can dominate the sum and push $\rho$ close to 1 while
$\delta_m$ remains poorly matched.
This is precisely what is observed: modes with $\rho = 0.999$ show a systematic
offset in $\delta_m$ not present in the other three fields.

\paragraph{Problem 2: correlation vs.\ equality.}
Even if the per-field contribution of $\delta_m$ to $\rho$ were 1, this would
only guarantee that $\Xn{1}{\delta_m}\va$ and $\Xn{2}{\delta_m}\vb$ are
\emph{proportional}, i.e.\
$\Xn{1}{\delta_m}\va = \lambda\,\Xn{2}{\delta_m}\vb$ for some scalar $\lambda$.
The Frobenius pre-normalisation further obscures the relationship between the
unnormalised time series.
Correlation is scale-invariant; equality is not.

%----------------------------------------------------------------------
\section{The Matching Solution Formulation}
\label{sec:matching}
%----------------------------------------------------------------------

\subsection{Residual objective}
\label{sec:residual}

Define the per-field residual vector for coefficient pair $(\va, \vb)$:
\begin{equation}
    \bm{r}^p(\va,\vb) = \Xn{1}{p}\va - \Xn{2}{p}\vb.
    \label{eq:residual}
\end{equation}
The goal~\eqref{eq:goal} (after Frobenius pre-normalisation) is $\bm{r}^p = \bm{0}$
for all $p$.
We seek non-trivial $(\va, \vb)$ that minimise the total squared residual:
\begin{equation}
    J(\va,\vb) = \sum_p w_p \lVert\bm{r}^p(\va,\vb)\rVert^2
               = \sum_p w_p \lVert\Xn{1}{p}\va - \Xn{2}{p}\vb\rVert^2.
    \label{eq:objective}
\end{equation}
Unlike the CCA correlation, $J = 0$ guarantees~\eqref{eq:goal} exactly.
The value of $J$ provides an absolute, per-field measure of disagreement:
it is not normalised by any self-energy, so a small $J$ cannot be achieved
by making $\va$ or $\vb$ small unless they are truly zero.

\subsection{Block matrix construction}
\label{sec:block}

Expand the squared norm in~\eqref{eq:objective}:
\begin{equation}
    \lVert\Xn{1}{p}\va - \Xn{2}{p}\vb\rVert^2
    = \va^\top \Gb{1}^p \va - 2\,\va^\top \Gb{12}^p \vb + \vb^\top \Gb{2}^p \vb,
\end{equation}
where the \emph{per-field} Gram matrices are
\begin{equation}
    \Gb{1}^p = (\Xn{1}{p})^\top\Xn{1}{p}, \quad
    \Gb{2}^p = (\Xn{2}{p})^\top\Xn{2}{p}, \quad
    \Gb{12}^p = (\Xn{1}{p})^\top\Xn{2}{p},
    \label{eq:pergram}
\end{equation}
each of shape $(\Nk,\Nk)$.
Note these are the same objects used in \texttt{cca\_formulation.pdf}, but here
kept separate per field rather than summed.

Concatenate $\vc = (\va^\top, \vb^\top)^\top \in \mathbb{R}^{2\Nk}$.
The objective becomes a quadratic form:
\begin{equation}
    J(\vc) = \vc^\top M\,\vc,
    \label{eq:quadratic}
\end{equation}
where $M \in \mathbb{R}^{2\Nk\times 2\Nk}$ is the symmetric positive semi-definite
block matrix
\begin{equation}
    M = \sum_p w_p
    \begin{pmatrix}
        \Gb{1}^p & -\Gb{12}^p \\[4pt]
        -(\Gb{12}^p)^\top & \Gb{2}^p
    \end{pmatrix}.
    \label{eq:M}
\end{equation}
Note that $M$ is the weighted sum of matrices of the form
$(Z^p)^\top Z^p \succeq 0$ where $Z^p = [\Xn{1}{p},\,-\Xn{2}{p}]$,
so $M \succeq 0$ always.
The diagonal blocks of $M$ recover the aggregate Gram matrices
$\Gb{1} = \sum_p w_p \Gb{1}^p$ and $\Gb{2} = \sum_p w_p \Gb{2}^p$
from \texttt{cca\_formulation.pdf}; the off-diagonal block is
$-\Gb{12} = -\sum_p w_p \Gb{12}^p$.

\subsection{Eigenvalue problem}
\label{sec:eigen}

Minimising $\vc^\top M\,\vc$ subject to $\|\vc\|=1$ is the standard minimum
eigenvalue problem.
Since $M \succeq 0$, its eigenvalues $\lambda_1 \leq \lambda_2 \leq \cdots$
are non-negative, and the minimisers are the eigenvectors corresponding to the
smallest eigenvalues:
\begin{equation}
    M\,\vc_i = \lambda_i\,\vc_i.
    \label{eq:eig}
\end{equation}
Each eigenvector $\vc_i = (\va_i^\top, \vb_i^\top)^\top$ splits into the
basis-1 coefficient $\va_i = \vc_i[:N_k]$ and the basis-2 coefficient
$\vb_i = \vc_i[N_k:]$.

The eigenvalue $\lambda_i$ is the minimum achievable value of the objective:
\begin{equation}
    \lambda_i = J(\vc_i) = \sum_p w_p \lVert\Xn{1}{p}\va_i - \Xn{2}{p}\vb_i\rVert^2.
\end{equation}
A mode with $\lambda_i \approx 0$ is a valid joint solution for \emph{all}
perturbation types simultaneously.
Modes with $\lambda_i$ large are incompatible between the two bases.

To set a threshold, it is convenient to normalise by the maximum eigenvalue
and keep modes with
\begin{equation}
    \frac{\lambda_i}{\lambda_{\max}} < \lambda_{\min},
    \label{eq:threshold}
\end{equation}
for a chosen $\lambda_{\min}$ (e.g.\ $10^{-3}$).
Alternatively, the raw value of $\lambda_i$ can be inspected directly: since
the data are Frobenius-normalised, $\lambda_i = 0$ is exact agreement.

\subsection{Null space interpretation}
\label{sec:null}

The condition~\eqref{eq:goal} is equivalent to
\begin{equation}
    Z^p\,\vc = \bm{0} \qquad \text{for all } p,
    \qquad Z^p = [\Xn{1}{p},\,-\Xn{2}{p}] \in \mathbb{R}^{\Nt\times 2\Nk}.
    \label{eq:null}
\end{equation}
So the exact solutions lie in the intersection of null spaces:
\begin{equation}
    \vc \in \ker(Z^{\delta_r}) \cap \ker(Z^{\delta_m})
           \cap \ker(Z^{v_r}) \cap \ker(Z^{v_m}).
    \label{eq:intersection}
\end{equation}
In general this intersection may be trivial (containing only $\bm{0}$), in which
case no exact joint eigenfunction exists.
The minimum-eigenvalue problem finds the vector $\vc$ that is ``closest'' to
being in all null spaces simultaneously, in the least-squares sense.

%----------------------------------------------------------------------
\section{Relationship to CCA}
\label{sec:relation}
%----------------------------------------------------------------------

The CCA and the matching-solution method optimise different objectives:

\begin{center}
\begin{tabular}{lll}
\toprule
 & CCA & Matching solution \\
\midrule
\textbf{Objective} & Maximise $\rho(\va,\vb) \in [0,1]$ (correlation)
                   & Minimise $J(\va,\vb) \geq 0$ (squared difference) \\
\textbf{Found when} & $\va \propto \vb$ (same shape)
                    & $\Xb{1}{p}\va = \Xb{2}{p}\vb$ (same value) \\
\textbf{Quality metric} & $\rho \to 1$ (high = good)
                        & $\lambda \to 0$ (low = good) \\
\textbf{Outvoting} & Yes: 3 fields can hide $\delta_m$ mismatch
                   & No: $J$ penalises every field independently \\
\textbf{Scale sensitivity} & No (correlation is scale-invariant)
                            & Yes ($\lambda = 0$ requires exact equality) \\
\textbf{Matrix size} & $(\Nk\times\Nk)$ Gram matrices
                     & $(2\Nk\times 2\Nk)$ block matrix \\
\bottomrule
\end{tabular}
\end{center}

\begin{theorem}
\label{thm:equiv}
If $\Xb{1}{p}\vc_0 = \Xb{2}{p}\vc_0$ for all $p$ for some $\vc_0 \neq \bm{0}$,
then $\vc = (\vc_0^\top, \vc_0^\top)^\top / \|\vc_0\|$ achieves $\lambda = 0$
in the matching-solution method and $\rho = 1$ in CCA (when
$\frob{\Xb{1}{p}} = \frob{\Xb{2}{p}}$ for all $p$).
\end{theorem}

\begin{proof}
Direct substitution: $J(\vc) = \sum_p w_p\|\Xn{1}{p}\vc_0/\|\vc_0\|
- \Xn{2}{p}\vc_0/\|\vc_0\|\|^2 = 0$, so $\lambda = 0$.
The CCA claim follows from Theorem~1 of \texttt{cca\_formulation.pdf}.
\end{proof}

The converse does not hold for CCA: $\rho = 1$ implies only proportionality,
not equality.
The converse \emph{does} hold for the matching-solution method: $\lambda = 0$
and $\vc\neq\bm{0}$ imply $\Xn{1}{p}\va = \Xn{2}{p}\vb$ for all $p$.

%----------------------------------------------------------------------
\section{Sparse Basis via Oblique Rotation}
\label{sec:sparsity}
%----------------------------------------------------------------------

The eigenvectors $\vc_i$ are dense in general.
Let $A = [\va_1 \mid \cdots \mid \va_m]$ and
$B = [\vb_1 \mid \cdots \mid \vb_m]$ collect the $m$ valid coefficient vectors.
We seek an invertible $M \in \mathbb{R}^{m\times m}$ such that the rotated
columns of $\tilde{A} = AM^\top$ are sparse (concentrated on one or a few
$k$-modes), applying the same rotation to $B$.

The same analysis as in \texttt{cca\_formulation.pdf} applies: Varimax is
inappropriate (the $\va_i$ are not orthonormal in $k$-space), while Promax
oblique rotation or FastICA are recommended.

%----------------------------------------------------------------------
\section{Complete Pipeline}
\label{sec:pipeline}
%----------------------------------------------------------------------

\begin{algorithm}[H]
\caption{Matching-solution sparse joint eigenfunction method}
\label{alg:pipeline}
\begin{algorithmic}[1]
\Statex \textbf{Input:} $\Xb{1}{p}, \Xb{2}{p} \in \mathbb{R}^{\Nt\times\Nk}$
        for each $p$; weights $\{w_p\}$; threshold $\lambda_{\min}$

\State \textbf{Frobenius normalise}
\State \quad $\Xn{\alpha}{p} \gets \Xb{\alpha}{p}/\frob{\Xb{\alpha}{p}}$
       for all $\alpha, p$

\State \textbf{Compute per-field Gram matrices} \Comment{Eq.~\eqref{eq:pergram}}
\State \quad $\Gb{1}^p \gets (\Xn{1}{p})^\top\Xn{1}{p}$ for each $p$
\State \quad $\Gb{2}^p \gets (\Xn{2}{p})^\top\Xn{2}{p}$
\State \quad $\Gb{12}^p \gets (\Xn{1}{p})^\top\Xn{2}{p}$

\State \textbf{Construct block matrix} \Comment{Eq.~\eqref{eq:M}}
\State \quad $M \gets \sum_p w_p
        \bigl[\begin{smallmatrix} \Gb{1}^p & -\Gb{12}^p \\
        -(\Gb{12}^p)^\top & \Gb{2}^p \end{smallmatrix}\bigr]$
       \Comment{$(2\Nk, 2\Nk)$}

\State \textbf{Solve eigenvalue problem} \Comment{Eq.~\eqref{eq:eig}}
\State \quad $M\,\vc_i = \lambda_i\,\vc_i$, \quad
       $\lambda_1 \leq \lambda_2 \leq \cdots$
\State \quad Keep $m$ eigenvectors with $\lambda_i/\lambda_{\max} < \lambda_{\min}$
\State \quad $\va_i \gets \vc_i[:N_k]$, \quad
       $\vb_i \gets \vc_i[N_k:]$ for $i = 1,\ldots,m$

\State \textbf{Oblique rotation for sparsity} \Comment{\S\ref{sec:sparsity}}
\State \quad $A \gets [\va_1\mid\cdots\mid\va_m]$,\quad
       $B \gets [\vb_1\mid\cdots\mid\vb_m]$
\State \quad $\tilde{A}, \tilde{B} \gets \mathrm{Promax}(A, B)$
       \Comment{or FastICA}

\Statex \textbf{Output:} $\tilde{A}, \tilde{B} \in \mathbb{R}^{\Nk\times m}$
\Statex \quad Each column pair $(\tilde{\va}_j, \tilde{\vb}_j)$: sparse coefficients
        satisfying $\Xb{1}{p}\tilde{\va}_j \approx \Xb{2}{p}\tilde{\vb}_j$
        for all $p$
\end{algorithmic}
\end{algorithm}

%----------------------------------------------------------------------
\section{Discussion}
\label{sec:discussion}
%----------------------------------------------------------------------

\begin{enumerate}
    \item \textbf{Stricter criterion.} The eigenvalue $\lambda_i$ measures
    the total squared difference across all perturbation types.
    No field can be ``outvoted'': even a small mismatch in a single field
    contributes positively to $\lambda_i$.
    This directly addresses the $\delta_m$ mismatch observed with the CCA method.

    \item \textbf{Larger matrix.} The block matrix $M$ is $(2\Nk\times 2\Nk)$
    compared to $(\Nk\times\Nk)$ for the CCA Gram matrices.
    For $\Nk \lesssim 10^3$ this remains computationally cheap; only the
    smallest few eigenvalues need to be computed (e.g.\ via
    \texttt{scipy.sparse.linalg.eigsh}).

    \item \textbf{Separate coefficient vectors.} The method naturally produces
    $\va_i \neq \vb_i$ in general.
    For exact joint eigenfunctions the two coincide; for approximate ones
    ($\lambda_i$ small but non-zero) they may differ slightly, and both are
    meaningful.

    \item \textbf{Weighting.} The weights $w_p$ encode the relative importance
    of each perturbation type.
    Unlike in CCA, a larger weight for $\delta_m$ now directly penalises
    $\delta_m$ mismatch rather than merely up-weighting its contribution to an
    aggregate correlation.
    Setting all $w_p = 1$ gives equal treatment to all fields.

    \item \textbf{Relation to CCA.} The CCA method can be recovered
    approximately from the matching-solution method.
    When the two bases are similar ($\frob{\Xb{1}{p}} \approx \frob{\Xb{2}{p}}$),
    modes with $\lambda_i \approx 0$ also achieve $\rho_i \approx 1$, but
    the converse does not hold.
    Use CCA when shape agreement is sufficient; use the matching-solution method
    when exact pointwise equality of time evolutions is required.
\end{enumerate}

\end{document}
